\chapter{Light-Ray Operators and Energy Correlators}
\label{ch:cft-light-ray-energy-correlators}

\ConventionsMixedCFT

Energy correlators provide a common observable language for conformal field
theory and high-energy gauge theory.  The primitive object is a detector at
future null infinity, represented in local QFT by an asymptotic flux of the
stress tensor.  The same object can also be described, after conformal
compactification, as a null-integrated light-ray operator.

\section{Detector States and Energy Flux}

Let the Lorentzian CFT live on \(\mathbb M^D\), \(D\geq3\), with Hilbert space
\(\Hilb\), vacuum vector \(\vac\), translation generators \(P^\mu\), and a
conserved symmetric stress tensor \(T^{\mu\nu}\).  A normalizable collider
state is obtained by smearing a local operator with timelike momentum support.
For a local operator \(\mathcal O_A\), where \(A\) denotes all tensor and
internal indices, choose a smooth compactly supported wavepacket
\(h(q)\) on the open future cone and a polarization tensor \(\varepsilon^A\),
and set
\[
  \ket{\Psi_{\mathcal O,h,\varepsilon}}
  =
  \int_{V_+}\dd\mu(q)\,
  h(q)\,\varepsilon^A\,\widetilde{\mathcal O}_A(q)\ket{\vac}.
\]
Here \(\dd\mu(q)\) is a positive measure on the chosen momentum slice.  The
Fourier-transformed operator is an operator-valued distribution applied to the
test function \(h\).  Sharp momentum states used in collider formulas are
distributional limits of these wavepackets.

\begin{hypothesis}[CFT energy detector]
\label{hyp:cft-energy-detector}
For every real \(f\in C^\infty(S^{D-2})\), the limit
\begin{equation}
  \mathcal E(f)
  =
  \mathrm{w}\!-\!\lim_{r\to\infty}
  r^{D-2}
  \int_{-\infty}^{\infty}\dd u
  \int_{S^{D-2}}\dd\Omega(\mathbf n)\,
  f(\mathbf n)\,
  n_i T^{0i}(u+r,r\mathbf n)
  \label{eq:cft-energy-detector}
\end{equation}
exists as a quadratic form on the collider domain generated by such states.
If \(f\geq0\), the resulting quadratic form is positive.  The constant
function gives the total outgoing energy,
\[
  \mathcal E(1)=P^0
\]
on the same domain.
\end{hypothesis}

The distributional detector density is denoted by
\[
  \mathcal E(f)
  =
  \int_{S^{D-2}}\dd\Omega(\mathbf n)\,
  f(\mathbf n)\mathcal E(\mathbf n).
\]
For \(k\) test functions \(f_1,\ldots,f_k\), the \(k\)-point energy correlator
in a normalized state \(\Psi\) is
\begin{equation}
  \mathcal G_k(f_1,\ldots,f_k;\Psi)
  =
  \bra{\Psi}\mathcal E(f_1)\cdots\mathcal E(f_k)\ket{\Psi}.
  \label{eq:cft-energy-correlator}
\end{equation}
The product is defined first for separated angular supports and then extended,
when possible, as a distribution on \((S^{D-2})^k\).  Coincident detector
directions carry contact terms determined by the distributional product of
light-ray operators.

\section{Relation to Null-Integrated Operators}

Let \(n^\mu=(1,\mathbf n)\) be a future null vector.  For a point
\(x_\perp\) in a spacelike plane transverse to \(n^\mu\), the averaged null
energy operator on the null line through \(x_\perp\) is the quadratic form
\begin{equation}
  \mathcal A_n(x_\perp)
  =
  \int_{-\infty}^{\infty}\dd\lambda\,
  T_{\mu\nu}(x_\perp+\lambda n)n^\mu n^\nu .
  \label{eq:anec-light-ray-operator}
\end{equation}
The conformal map from a null line to future null infinity identifies
\(\mathcal E(\mathbf n)\), up to the fixed conformal Jacobian determined by
the chosen compactification, with the light-ray transform
\eqref{eq:anec-light-ray-operator} of the stress tensor.  Thus the detector
operator is a nonlocal operator obtained by integrating a local field along a
complete null generator.

This formulation separates three pieces of data:
\[
  \begin{gathered}
  \text{local stress tensor}
  \longrightarrow
  \text{null-integrated light-ray operator}
  \\
  \longrightarrow
  \text{calorimetric distribution on }S^{D-2}.
  \end{gathered}
\]
The first belongs to the local operator algebra.  The second requires control
of a null integral of an operator-valued distribution.  The third is a
detector observable in states with finite total energy.

For \(f\geq0\), positivity of \(\mathcal E(f)\) is the detector version of the
average null energy condition.  In a unitary Lorentzian CFT satisfying the
hypotheses of the ANEC theorem, the quadratic form
\(\int\dd\lambda\,T_{\mu\nu}(x_\perp+\lambda n)n^\mu n^\nu\) is positive.
In this chapter that positivity is recorded in
Hypothesis~\ref{hyp:cft-energy-detector}, because the subsequent collider
inequalities use the detector as an operator on states.

\section{Conformal Collider One-Point Functions}

Consider a four-dimensional parity-even CFT and a stress-tensor state in the
rest frame \(q^\mu=(Q,\mathbf0)\).  Let \(\varepsilon_{ij}\) be a complex
symmetric traceless spatial polarization, \(i,j=1,2,3\), and write
\[
  \ket{\Psi_{\varepsilon,q}}
  =
  \varepsilon_{ij}\widetilde T^{ij}(q)\ket{\vac}
\]
with wavepacket regularization understood.  Rotational covariance around
\(\mathbf q=\mathbf0\), stress-tensor conservation, and the Ward identity
normalization of \(T_{\mu\nu}\) imply that the separated-angle one-detector
expectation has the form
\begin{equation}
  \frac{\bra{\Psi_{\varepsilon,q}}\mathcal E(\mathbf n)
  \ket{\Psi_{\varepsilon,q}}}
  {\braket{\Psi_{\varepsilon,q}}{\Psi_{\varepsilon,q}}}
  =
  \frac{Q}{4\pi}
  \left[
    1
    +t_2
    \left(
      \frac{\varepsilon^*_{ij}\varepsilon_{ik}n^jn^k}
           {\varepsilon^*_{ij}\varepsilon_{ij}}
      -\frac13
    \right)
    +t_4
    \left(
      \frac{|\varepsilon_{ij}n^in^j|^2}
           {\varepsilon^*_{ij}\varepsilon_{ij}}
      -\frac{2}{15}
    \right)
  \right].
  \label{eq:hm-energy-one-point}
\end{equation}
The real constants \(t_2,t_4\) are coordinates on the parity-even
stress-tensor three-point function after \(C_T\) and Ward-identity contact
terms have been fixed.

\begin{proposition}[Four-dimensional conformal-collider inequalities]
\label{prop:hm-inequalities}
Assume the energy detector in
\eqref{eq:hm-energy-one-point} is a positive quadratic form for every
direction \(\mathbf n\).  Then
\[
  1-\frac13t_2-\frac{2}{15}t_4\geq0,
  \qquad
  1+\frac16t_2-\frac{2}{15}t_4\geq0,
  \qquad
  1+\frac13t_2+\frac{8}{15}t_4\geq0.
\]
\end{proposition}

\begin{proof}
Fix \(\mathbf n=(0,0,1)\).  The little group \(SO(2)\) decomposes symmetric
traceless polarizations into helicity \(2\), helicity \(1\), and helicity
\(0\) sectors.  For a helicity \(2\) polarization, for instance
\(\varepsilon_{11}=-\varepsilon_{22}\) and
\(\varepsilon_{12}=0=\varepsilon_{i3}\), both ratios in
\eqref{eq:hm-energy-one-point} equal \(0\).  Positivity gives
\[
  1-\frac13t_2-\frac{2}{15}t_4\geq0.
\]
For a helicity \(1\) polarization, for instance
\(\varepsilon_{13}=\varepsilon_{31}\neq0\), the first ratio equals
\(\frac12\) and the second equals \(0\), giving the second inequality.  For a
helicity \(0\) polarization proportional to
\(\operatorname{diag}(1,1,-2)\), the first and second ratios both equal
\(\frac23\), giving the third inequality.  Rotational covariance transports
the same conditions to every detector direction.
\end{proof}

The derivation of the normal form
\eqref{eq:hm-energy-one-point} is a calculation of a Wightman
\(\langle TTT\rangle\) distribution integrated against the detector kernel.
Contact terms in the local \(\langle TTT\rangle\) Ward identities fix the
normalization of the stress tensor and do not change separated-angle detector
values after the state wavepacket and detector smearing have been specified.

\section{Energy-Energy Correlators in a CFT}

For a state of sharp total energy \(Q\), the normalized energy-energy
correlator is the distribution
\begin{equation}
  \mathrm{EEC}_\Psi(\chi)
  =
  \frac{1}{Q^2\braket{\Psi}{\Psi}}
  \bra{\Psi}
  \int_{S^{D-2}}\dd\Omega_1\dd\Omega_2\,
  \delta(\cos\chi-\mathbf n_1\cdot\mathbf n_2)
  \mathcal E(\mathbf n_1)\mathcal E(\mathbf n_2)
  \ket{\Psi}.
  \label{eq:cft-eec}
\end{equation}
For a wavepacket in \(Q\), replace \(Q^2\) by the corresponding expectation
of \((P^0)^2\), or keep the unnormalized numerator and denominator separately.
The separated-angle CFT observable
\eqref{eq:cft-eec} is the same type of detector distribution as the QCD
observable \eqref{eq:qcd-eec-definition}.  The difference is dynamical:
in a CFT there is no mass gap and no hadronic scattering basis, while the
detector remains defined by the stress tensor and null infinity.

When two detector directions approach one another, the product of light-ray
operators has a short-distance expansion on the celestial sphere.  The
following theorem is used as an external theorem boundary for the modern
light-ray formalism.\footnote{See M.~Kologlu, P.~Kravchuk,
D.~Simmons-Duffin, and A.~Zhiboedov, ``The light-ray OPE and conformal
colliders,'' arXiv:1905.01311.  For the conformal-collider formulation, see
D.~M.~Hofman and J.~Maldacena, ``Conformal collider physics: energy and charge
correlations,'' JHEP 05 (2008) 012, arXiv:0803.1467.}

\begin{theorem}[Light-ray OPE under Lorentzian CFT hypotheses]
\label{thm:light-ray-ope-boundary}
In a unitary CFT satisfying the Lorentzian analyticity, spectrum, and
operator-growth hypotheses required for the light-ray OPE, products of
null-integrated operators on the same null plane admit a convergent expansion
in light-ray operators.  For energy detectors, the separated-angle product has
the form, as a distribution tested near \(\mathbf n_1=\mathbf n_2\),
\[
  \mathcal E(\mathbf n_1)\mathcal E(\mathbf n_2)
  =
  \sum_\alpha
  C_\alpha(\mathbf n_1,\mathbf n_2)\,
  \mathbb L_\alpha(\mathbf n_2),
\]
where \(\mathbb L_\alpha\) are light-ray transforms of local CFT primaries and
their descendants, analytically continued in spin as specified by the
Lorentzian inversion formula.
\end{theorem}

This theorem gives a nonperturbative meaning to the collinear expansion of
energy correlators in a CFT.  It is the conformal counterpart of the QCD
statement that small-angle energy correlators are governed by collinear
operator data and anomalous dimensions.  In QCD the running coupling and the
mass gap alter the long-distance behavior; in a CFT the light-ray OPE is an
operator expansion intrinsic to the fixed point.

\section{Comparison with QCD Detector Observables}

The QCD and CFT constructions share the same operator-theoretic skeleton:
\[
  T^{\mu\nu}
  \quad\leadsto\quad
  \mathcal E(\mathbf n)
  \quad\leadsto\quad
  \mathcal G_k(\mathbf n_1,\ldots,\mathbf n_k).
\]
In QCD, the physical final-state resolution is a hadronic spectral sum and
perturbation theory enters through asymptotically free short-distance
coefficient functions.  In a CFT, the state can be created directly by local
operators and the angular dependence is constrained by conformal
representation theory.  The common lesson is mathematical rather than
phenomenological: a detector observable is defined by the stress tensor in the
physical Hilbert space, and partonic or conformal expansions are calculational
representations of that observable under stated hypotheses.
